% ============================================================
% Section 4 — Experimental Setup
% ============================================================
\section{Experiments}
\label{sec:experiments}

\subsection{Benchmarks and Tasks}
\label{sec:experiments:benchmarks}

We evaluate HAT on three families of tasks that probe different aspects of self-improvement:

\paragraph{Knowledge-intensive question answering.}
We use Natural Questions (NQ)~\citep{kwiatkowski2019natural} and TriviaQA~\citep{joshi2017triviaqa} in a \emph{streaming} setup: questions arrive sequentially and the model must answer without access to a retrieval corpus.  After each answer, the model may receive a binary correctness signal or an explicit user correction.

\paragraph{Instruction following.}
We adopt a subset of the Alpaca-Eval benchmark~\citep{li2023alpacaeval} and MT-Bench~\citep{zheng2023judging}, where instructions are presented in temporal batches that gradually increase in complexity.  This tests HAT's ability to internalise and generalise from corrective feedback on instruction-following quality.

\paragraph{Personalisation.}
We construct a personalisation benchmark based on LaMP~\citep{salemi2024lamp}, where a simulated user provides preferences and corrections over successive interactions.  The model is evaluated on how well it adapts to individual user style and factual preferences over time.

\subsection{Baselines}
\label{sec:experiments:baselines}

We compare HAT against the following approaches:
\begin{itemize}[leftmargin=*,itemsep=2pt]
  \item \textbf{Frozen}: the base Cortex model without any adaptation.
  \item \textbf{RAG}~\citep{lewis2020retrieval}: retrieval-augmented generation using a dense retriever over the interaction history.
  \item \textbf{Na\"ive Replay}: continual fine-tuning on all interaction logs without selection.
  \item \textbf{EWC Replay}~\citep{kirkpatrick2017overcoming}: replay with elastic weight consolidation but without experience curation.
  \item \textbf{DER++}~\citep{buzzega2020dark}: dark experience replay with knowledge distillation, a strong continual learning baseline.
  \item \textbf{Self-Play}~\citep{chen2024self}: the model generates and filters its own training data via self-consistency.
  \item \textbf{Full Distillation}: fine-tuning the student on the oracle's responses for \emph{all} queries (upper bound on oracle usage).
\end{itemize}

\subsection{Implementation Details}
\label{sec:experiments:details}

\paragraph{Base models.}
We instantiate the Cortex with three model families at different scales: Qwen-2.5-1.5B, Phi-3-mini-3.8B, and LLaMA-3-8B.  All models are instruction-tuned variants.

\paragraph{Oracle.}
We use GPT-4o as the default oracle $\oracle$.  We also report results with a smaller oracle (LLaMA-3-70B) to study oracle quality sensitivity.

\paragraph{Hippocampus Agent.}
The abstraction module uses the Cortex itself with dedicated triage and extraction prompts, making HAT lightweight and self-contained.  Write eligibility is determined by the log-probability uncertainty gate in Eq.~\eqref{eq:uncertainty_gate}.  Semantic redundancy is handled separately by the embedding-based CREATE/REVISE router in Eq.~\eqref{eq:create_revise}, rather than being mixed into the write score.

\paragraph{Hyperparameters.}
We set the write threshold to $\tau_w = 0.3$ in Eq.~\eqref{eq:write_decision} and the memory-routing threshold to $\tau_r = 0.85$ in Eq.~\eqref{eq:create_revise}, based on validation-set tuning.  Oracle consultation is triggered when $\uncertainty > \tau_\oracle = 0.7$.  SWS training uses LoRA (rank 16) with AdamW ($\eta = 2 \times 10^{-5}$), batch size 32, and 3 epochs per cycle.  The EWC coefficient is $\mu = 0.1$.  A SWS cycle is triggered every 500 interactions.

\paragraph{Evaluation protocol.}
We report metrics averaged over 3 random seeds.  Each experiment simulates 5{,}000 sequential interactions (10 SWS cycles).  We measure: (1) \textbf{accuracy} (exact match or ROUGE-L depending on task), (2) \textbf{error rate} (fraction of incorrect responses), (3) \textbf{sample efficiency} (accuracy gain per training example), and (4) \textbf{forgetting rate} (performance drop on earlier tasks after adaptation).
